\documentclass[11pt]{article}

\usepackage[T1]{fontenc}
\usepackage[utf8]{inputenc}
\usepackage[margin=1in]{geometry}
\usepackage{amsmath,amssymb,amsthm,mathtools,stmaryrd}
\usepackage{mathpartir}
\usepackage{xcolor}
\usepackage{graphicx}
\usepackage{enumitem}
\usepackage{hyperref}
\hypersetup{colorlinks=true,linkcolor=black,citecolor=black,urlcolor=blue}

% ---- colour discipline (red copy vs black copy) -------------------------
\definecolor{redcopy}{rgb}{0.78,0.07,0.07}
\newcommand{\rd}[1]{\textcolor{redcopy}{#1}}

% ---- four sorts --------------------------------------------------------
\newcommand{\RS}{\rd{\mathbf{S}}}              % red set
\newcommand{\RA}{\rd{\mathbf{A}}}              % red atom
\newcommand{\BS}{\mathbf{S}}                   % black set
\newcommand{\BA}{\mathbf{A}}                   % black atom

% ---- structural maps ----------------------------------------------------
\newcommand{\swp}{\mathsf{s}}                  % colour swap (involution)
\newcommand{\kt}{\kappa}                       % red knot  : BlackSet -> RedAtom
\newcommand{\ktb}{\bar\kappa}                  % black knot : RedSet  -> BlackAtom

% ---- process syntax -----------------------------------------------------
\newcommand{\Zero}{\mathbf{0}}
\newcommand{\Par}{\mathbin{\mid}}
\newcommand{\inp}[3]{\mathsf{for}(#1 \leftarrow #2)\,#3}
\newcommand{\out}[2]{#1\mathbin{!}(#2)}
\newcommand{\drp}[1]{{\ast}#1}
\newcommand{\quo}[1]{@#1}
\newcommand{\abs}[2]{(#1)\,#2}                  % abstraction (y)P
\newcommand{\conc}[1]{[#1]}                     % concretion [Q]
\newcommand{\nset}[1]{[\mathcal{N}]#1}

% ---- semantic apparatus -------------------------------------------------
\newcommand{\Proc}{\mathsf{Proc}}
\newcommand{\Name}{\mathsf{Name}}
\newcommand{\Set}{\mathrm{Set}}
\newcommand{\Pfin}{\mathcal{P}_{\mathrm{fin}}}
\newcommand{\B}{\mathsf{B}}
\newcommand{\sem}[1]{\llparenthesis #1\rrparenthesis}   % behavioural denotation
\newcommand{\lts}[1]{\xrightarrow{\,#1\,}}
\newcommand{\Sym}{\mathrm{Sym}}
\newcommand{\supp}{\mathrm{supp}}
\newcommand{\bis}{\sim}
\newcommand{\fresh}{\mathbin{\#}}
\newcommand{\scong}{\equiv}
\newcommand{\nameq}{\equiv_{\mathcal N}}
\newcommand{\rpar}{\mathbin{\|}}               % semantic parallel
\newcommand{\react}{\rightarrowtail}           % reaction
\newcommand{\hole}{[\,\cdot\,]}                % context hole
\newcommand{\DK}{\partial T_{K}}               % Fire context-label set
\newcommand{\plug}[2]{#1[#2]}                  % plug

\theoremstyle{plain}
\newtheorem{theorem}{Theorem}
\newtheorem{proposition}[theorem]{Proposition}
\newtheorem{lemma}[theorem]{Lemma}
\newtheorem{corollary}[theorem]{Corollary}
\theoremstyle{definition}
\newtheorem{definition}[theorem]{Definition}
\newtheorem{remark}[theorem]{Remark}
\newtheorem{obligation}{Obligation}

\title{\bf Quoting is Colour-Swap:\\ a model of the rho calculus in the knotted universe}
\author{Lucius Gregory Meredith\\[2pt]
\small CEO, F1R3FLY.io, and Chief Science Officer, F1R3FLY Industries\\
\small \texttt{lgreg.meredith@gmail.com}}
\date{}

\begin{document}
\maketitle

\begin{abstract}
The knotted universe of \emph{A Knotted Universe: a new notion of reflective set
theories} closes with a conjecture: that the reflective higher-order (rho)
calculus has a natural set-theoretic model there, with the quote $\quo{\,}$ read
as sealing a process into an atom of the dual colour and the dereference
$\drp{\,}$ as unsealing it. We make the conjecture good, and in a form that puts
the colour barrier itself to work. We \emph{stratify} the universe into four
sorts --- red sets and red atoms, black sets and black atoms --- linked by the two
knots and by a single polymorphic involution, the colour-swap $\swp$. The
calculus's reflection operators are then \emph{derived}: $\quo{\,}=\kt\circ\swp$
and $\drp{\,}=\swp\circ\kt^{-1}$, so that $\drp{\quo P}\scong P$ becomes
$\swp\circ\swp=\mathrm{id}$ --- the reflection law is the involutivity of the
swap, not an axiom about the model. We give the dynamics following Milner's
guidance for the polyadic $\pi$-calculus: a process denotes the set of its
\emph{commitments}, input committing to an abstraction $\abs y P$ and output to a
concretion $\conc Q$, with communication the application $\abs y P\cdot\conc M =
P\{\quo M/y\}$ --- and it is exactly here, in the reaction rule, that the swap
fires, sealing the sent process into the received name. The denotation lands in
the final coalgebra of a behaviour functor over the knotted nominal universe,
where the colour-blind non-well-foundedness supplies recursion for free and
finality makes bisimilar processes equal. Following the GSLT discipline we label
transitions not by primitive actions but by \emph{contexts} --- the minimal
environments enabling a reaction, after Leifer and Milner --- and express
bisimulation over these context labels; the labels are the program/environment
cuts, simultaneously the indices of the OSLF-generated modal logic. We prove a
\emph{full abstraction} theorem, $\sem P=\sem Q$ iff $P\bis Q$, by identifying
context bisimilarity with the kernel of the unique morphism to the final coalgebra;
because the labels are contexts, the resulting equivalence is a congruence by
construction, so full abstraction holds for a congruence with no closure step. The
risk ledger of the source paper transfers: finite processes cost no infinite risk;
recursion --- which in rho enters through reflection, not a primitive replication
--- costs one $\omega$, black infinity, and not one primitive name. Aczel's CCS
model is recovered as the equivariant, colour-respecting fragment.
\end{abstract}

\section{Introduction}

Aczel modelled Milner's CCS inside non-well-founded set theory \cite{aczel}: a
process $P$ is identified with the set of its labelled transitions
$\{(a,P'):P\lts a P'\}$, a looping process is a non-well-founded set, bisimilar
processes are \emph{equal}, and the anti-foundation axiom's solution lemma
individuates processes as the unique solutions of their behavioural equations.
The construction is canonical, but it buys two things by fiat: an external,
unstructured supply of action names $a$, and the non-well-foundedness itself.

The reflective higher-order calculus --- the rho calculus \cite{rho} --- differs
from CCS and the $\pi$-calculus precisely where Aczel spends his first postulate.
Rho has \emph{no} primitive names. A name is a quoted process: the quote
$\quo{\,}$ turns a process into a name, the dereference $\drp{\,}$ turns a name
back into the process it quotes, and the two compose to the identity,
$\drp{\quo P}\scong P$. There is correspondingly no restriction operator $\nu$ and
no scope extrusion: fresh names are \emph{computed}, by quoting.

The knotted universe of \cite{knot} was built from the empty set alone, as two
mutually recursive copies of set theory --- a red copy $\Set[X]$ and a black copy
$\Set[\bar X]$ --- tied by $X=\Set[\bar X]$, $\bar X=\Set[X]$, so that each
colour's \emph{atoms} are the other colour's \emph{sets}. Section~9.3 of
\cite{knot} conjectures this universe is the natural home for an Aczel-style model
of rho, with $\quo{\,}$ sealing a process into an atom of the opposite colour and
$\drp{\,}$ unsealing it. We carry the conjecture out. Three design decisions shape
the development.

\paragraph{Quoting is the colour-swap, read through a stratification (Section~\ref{sec:swap}).}
Rather than collapse ``red atom'' and ``black set'' into one sort, as the bare
knot does, we keep four sorts and make the colour-moving maps first class. A
single polymorphic involution $\swp$, the colour-swap, exchanges red and black
\emph{while preserving the set/atom distinction}; the two knots $\kt,\ktb$ supply
the set/atom coercions. The calculus's $\quo{\,}$ and $\drp{\,}$ are then
\emph{composites} of these --- $\quo{\,}=\kt\circ\swp$, $\drp{\,}=\swp\circ\kt^{-1}$
--- and the reflection law is $\swp\circ\swp=\mathrm{id}$. This is at once more
honest (the colour crossing is explicit, not hidden in a definitional identity)
and exactly the graded, colour-tracking discipline \cite{knot}, \S6.5 asks for:
the four sorts are the grading, and $\swp$ is the reflective move.

\paragraph{The dynamics follow Milner (Section~\ref{sec:denotation}).}
After Milner's polyadic $\pi$ tutorial \cite{polyadic}, a process denotes the set
of its immediate \emph{commitments}; input commits to an \emph{abstraction}
$\abs y P$ and output to a \emph{concretion} $\conc Q$, and communication is the
application $\abs y P\cdot\conc M = P\{\quo M/y\}$. The reaction rule is where the
swap of Section~\ref{sec:swap} does its dynamical work: it seals the sent process
$M$ into the name $\quo M$ the receiver binds. Parallel composition is not bare
union --- that would wrongly identify $\out a 0\Par\out a 0$ with $\out a 0$ ---
but the corecursive expansion $\rpar$ that keeps each component's surviving
context in the residuals.

\paragraph{Transitions are labelled by contexts (Sections~\ref{sec:lts},~\ref{sec:fa}).}
To stay consistent with the GSLT programme we do not label transitions by primitive
actions but by \emph{contexts}: $P\lts F P'$ means $F$ is a minimal environment ---
an idem-pushout, after Leifer and Milner \cite{leifermilner,sewell}, adapted to
GSLTs \cite{meredithwells} --- in which $P$ completes a reaction to $P'$. The label
is the program/environment cut at which $P$ meets its world, and it is at once the
Firing label of the transition system, the Witnessing index of an OSLF-generated
Hennessy--Milner modality $\langle F\rangle\varphi$, and the Placing of the
interaction in $P$'s term structure --- three views of one $\partial T$. Expressing
bisimulation over context labels has a structural dividend: the equivalence is a
congruence by construction, so the full-abstraction theorem holds for a congruence
with none of the closure-under-substitution care an action-labelled account
requires.

\paragraph{Recursion is free (Sections~\ref{sec:functor},~\ref{sec:fa}).}
The denotation lands in the final coalgebra of a behaviour functor over the knotted
nominal universe; the colour-blind non-well-foundedness of \cite{knot}, \S7
supplies the recursion (which in rho is reached through reflection, with no
primitive replication operator), and finality makes bisimilar processes equal. We
prove full abstraction, $\sem P=\sem Q\iff P\bis Q$, identifying context
bisimilarity with the kernel of the unique morphism to the final coalgebra, and
read it back through OSLF adequacy.

We work throughout with the modern presentation: $\quo P$ for quoting, $\drp x$
for dereference, $\inp y x P$ for input.

\section{The rho calculus}\label{sec:rho}

\subsection{Syntax}
The two-sorted grammar of processes $P$ and names $x$:
\[
  P,Q \;::=\; \Zero \;\mid\; P \Par Q \;\mid\; \inp y x P \;\mid\; \out x Q
            \;\mid\; \drp x,
  \qquad\qquad
  x,y \;::=\; \quo P .
\]
$\Zero$ is the stopped process; $P\Par Q$ parallel composition; $\inp y x P$ the
input receiving on channel $x$, binding the received name to $y$ in $P$;
$\out x Q$ the asynchronous output sending $Q$ on $x$; and $\drp x$ the
dereference. A name is solely a quoted process: no primitive names, and no
restriction $\nu$. In $\inp y x P$ the name $y$ is bound in $P$. Capture-avoiding
substitution $P\{x/y\}$ of a name for a free name reaches inside quotes.

There is no primitive replication or recursion. Both are nonetheless available, by
reflection: a recursive process is written by quoting its own continuation and
dereferencing it, e.g.\ a server
$D \scong \inp y x (\drp y \Par \out x y)\Par\out x{\quo{(\cdots)}}$ reinstalls
itself after each interaction. Infinite behaviour from a finite term is therefore
present already in the core calculus --- which is why the model needs the
non-well-founded reading even with no $\,!\,$ in the grammar.

\subsection{Structural congruence and name equivalence}
$\scong$ is the least congruence closed under $\alpha$-equivalence and
\[
  (P\Par Q)\Par R \scong P\Par(Q\Par R),\quad
  P\Par Q \scong Q\Par P,\quad
  P\Par\Zero \scong P,\quad
  \drp{\quo P}\scong P .
\]
Name equivalence is induced by dereference:
$\quo P \nameq \quo Q :\Longleftrightarrow \drp{\quo P}\scong\drp{\quo Q}
 \Longleftrightarrow P\scong Q$.

\subsection{Reaction and the context-labelled transition system}\label{sec:lts}
The base interaction is communication, with the structural propagation rules:
\begin{mathpar}
  \inferrule*[right=\textsc{Comm}]{x \nameq z}
            {\inp y x P \,\Par\, \out z Q \;\longrightarrow\; P\{\quo Q / y\}}
  \and
  \inferrule*[right=\textsc{Par}]{P\longrightarrow P'}{P\Par R\longrightarrow P'\Par R}
  \and
  \inferrule*[right=\textsc{Struct}]{P\scong P' \\ P'\longrightarrow Q' \\ Q'\scong Q}
            {P\longrightarrow Q}
\end{mathpar}
Rho is an interactive GSLT whose single base constructor is parallel composition:
$K=\Par$ is the head of the left-hand side of the one hypothesis-free rule,
\textsc{Comm}, and the cut it induces is \emph{relative} --- each parallel
component is, from the other's vantage, the environment.

To define bisimulation we follow the GSLT discipline and label transitions by
\emph{contexts} rather than by primitive actions, after the minimal-context
(relative-pushout) construction of Leifer and Milner \cite{leifermilner,sewell}, as
adapted to GSLTs \cite{meredithwells}. A \emph{context} $F$ is a process with one
hole $\hole$; $\plug F P$ is the result of plugging $P$. Communication never fires
under a prefix, so the reactive contexts are, modulo $\scong$, the parallel
contexts $\hole\Par R$: the cut always lies at a parallel composition, exactly the
program/environment cut.

\begin{definition}[Context-labelled transitions]\label{def:ctx-lts}
$P\lts F P'$ holds when $F$ is a context with $\plug F P\longrightarrow P'$ by a
single \textsc{Comm}, and $F$ is \emph{minimal} --- an idem-pushout, supplying just
the missing half of the redex and no more. Concretely, modulo $\scong$:
\begin{itemize}[nosep]
  \item \emph{(offer input)} if $P\scong\inp y x{P_0}\Par P_1$ then
        $P\lts{\ \hole\,\Par\,\out x Q\ } P_0\{\quo Q/y\}\Par P_1$, one transition
        per message $Q$;
  \item \emph{(offer output)} if $P\scong\out x Q\Par P_1$ then
        $P\lts{\ \hole\,\Par\,\inp y x R\ } P_1\Par R\{\quo Q/y\}$, one per input
        context $\inp y x R$;
  \item \emph{(internal redex)} if $P\scong\inp y x{P_0}\Par\out x Q\Par P_1$ then
        $P\lts{\ \hole\ } P_0\{\quo Q/y\}\Par P_1$.
\end{itemize}
The label records what environment $P$ needs to interact; the identity context
$\hole$ labels the reactions $P$ performs alone (the silent moves). The seal
$\quo{\,}=\kt\circ\swp$ acts in each substitution: communication delivers the
\emph{name} $\quo Q$ of the sent process, never the process itself.
\end{definition}

\begin{remark}[The label set is $\partial T$, in its three views]\label{rem:dT}
The labels are the Fire-sub-bundle $\DK$ of the splitting space $\partial T\times T$:
the splittings $(F,t)$ whose subterm $t$ completes a \textsc{Comm} redex. This one
set is read three ways \cite{meredithwells}: as the label set of the transition
system (\emph{Firing}); as the index set of the OSLF-generated Hennessy--Milner
modalities $\langle F\rangle\varphi$ --- ``some $F$-labelled transition reaches a
state satisfying $\varphi$'' (\emph{Witnessing}); and as the program/environment
cuts, i.e.\ the locations at which $P$ meets its environment (\emph{Placing}). We
will use the first to define bisimulation and the second to characterise it.
\end{remark}

\begin{definition}[Context bisimulation]\label{def:bis}
A symmetric relation $\mathcal R$ on processes is a \emph{context bisimulation} if
$P\mathrel{\mathcal R}Q$ implies: whenever $P\lts F P'$ there is $Q\lts F Q'$ with
$P'\mathrel{\mathcal R}Q'$ (and symmetrically). Context bisimilarity $\bis$ is the
largest context bisimulation; it contains $\scong$.
\end{definition}

\begin{proposition}[Congruence, by construction]\label{prop:cong}
If the GSLT minimal-context construction supplies relative pushouts for rho's
contexts \cite{leifermilner,sewell,meredithwells}, then context bisimilarity is a
congruence: $P\bis Q$ implies $\plug F P\bis\plug F Q$ for every context $F$.
\end{proposition}
\begin{proof}[Proof sketch]
This is the Leifer--Milner theorem: when reactions are presented as idem-pushout
labels over a category of contexts possessing relative pushouts, the derived
bisimilarity is a congruence, because the bisimulation game already quantifies over
the discriminating power of every context. The hypothesis --- that rho's contexts
have the requisite relative pushouts --- is the content of the GSLT minimal-context
construction \cite{meredithwells}, and is recorded as Obligation~\ref{ob:rpo}.
\end{proof}

This is the structural reason the GSLT programme labels by contexts: congruence is
\emph{built into} the equivalence rather than recovered afterwards. It settles, at
the level of the definition, the point that an action-labelled presentation must
negotiate by closing under name instantiation (compare Section~\ref{sec:fa}).

\begin{remark}[Commitments are the finite presentation]\label{rem:commit}
The offer-input clause of Definition~\ref{def:ctx-lts} is a family
$\{\,P\lts{\hole\Par\out x Q}P_0\{\quo Q/y\}\Par P_1 : Q\in\Proc\,\}$ indexed by all
messages. The family factors through a single \emph{abstraction} $\abs y{(P_0\Par P_1)}$
(with $y\fresh P_1$): plugging the message $\quo Q$ into the abstraction gives the
residual. Likewise the offer-output family factors through a \emph{concretion}.
These abstraction/concretion packagings, in Milner's sense \cite{polyadic}, are the
finite presentation of the class-indexed context-transition families; we use them
for the denotation of Section~\ref{sec:denotation}, with the understanding that the
labels they abbreviate are the contexts of Definition~\ref{def:ctx-lts}.
\end{remark}

\section{Quoting is the colour-swap: a stratified reading}\label{sec:swap}

\subsection{Four sorts and one involution}
The construction of \cite{knot} takes two copies of one signature-generated set
theory, red $\Set[X]$ and black $\Set[\bar X]$, and ties them by
$X=\Set[\bar X]$, $\bar X=\Set[X]$: each colour's atoms are the other's sets. The
source paper reads these ties as identities, collapsing ``red atom'' with ``black
set''. For the rho model it is cleaner --- and, as we will see, exactly the graded
discipline \cite{knot}, \S6.5 calls for --- to keep the sorts distinct and make
the colour-moving maps first class. We name four sorts,
\[
  \RS\ \text{(red sets)},\quad \RA\ \text{(red atoms)},\qquad
  \BS\ \text{(black sets)},\quad \BA\ \text{(black atoms)},
\]
and three families of maps.

\paragraph{The knots, as coercions.} The two recursive knots become two
isomorphisms,
\[
  \kt \colon \BS \xrightarrow{\ \cong\ } \RA
  \qquad(\text{black sets are red atoms; } X=\Set[\bar X]),
\]
\[
  \ktb \colon \RS \xrightarrow{\ \cong\ } \BA
  \qquad(\text{red sets are black atoms; } \bar X=\Set[X]).
\]
Where \cite{knot} writes $X=\Set[\bar X]$ as an equality we write $\kt$ as a
coercion: a black set, presented \emph{as} a red atom, is red-opaque
(\cite{knot}, Remark~2); $\kt$ is the act of so presenting it.

\paragraph{The swap, an involution.} The two copies are copies of one theory, so
the construction is symmetric under exchanging the colours --- the symmetry
\cite{knot} invokes for its Corollary~2. We make it a map: a single polymorphic,
\emph{sort-preserving} involution
\[
  \swp \colon \RS\cong\BS,\qquad \swp\colon\RA\cong\BA,
  \qquad \swp\circ\swp=\mathrm{id},
\]
repainting red to black and back, structure for structure ($\swp$ carries
$\rd\emptyset$ to $\emptyset$, $\rd\cup$ to $\cup$, and so on). The swap commutes
with the knots: the square
\[
  \swp\circ\kt \;=\; \ktb\circ\swp
  \qquad(\text{both } \BS \to \BA)
\]
holds, expressing that repainting a black-set-as-red-atom agrees with repainting
then re-presenting. This coherence is what lets one colour's reflection generate
the other's.

\subsection{The calculus's reflection, derived}
Take processes to be red sets and names to be red atoms. Define
\[
  \boxed{\;
    \quo{\,} \;:=\; \kt\circ\swp \;\colon\; \RS \to \RA,
    \qquad
    \drp{\,} \;:=\; \swp\circ\kt^{-1} \;\colon\; \RA \to \RS.
  \;}
\]
Thus $\quo P$ first repaints the red process $P$ black ($\swp$), then presents the
result as a red atom ($\kt$): it ``seals $P$ into an atom of the opposite colour''
exactly as \cite{knot}, \S9.3 proposes, the seal being the red-opacity of $\kt$,
the colour crossing being $\swp$. The dereference $\drp x$ un-presents
($\kt^{-1}$) and repaints back ($\swp$).

\begin{proposition}[Reflection is involutivity of the swap]\label{prop:reflection}
For every process $P$ and name $x$,
\[
  \drp{\quo P}
   = \swp\,\kt^{-1}\,\kt\,\swp\,(P)
   = \swp\,\swp\,(P) = P,
  \qquad
  \quo{\drp x}
   = \kt\,\swp\,\swp\,\kt^{-1}(x)
   = x.
\]
Hence the model validates $\drp{\quo P}\scong P$ and $\quo{\drp x}\nameq x$, not by
imposing an equation but as the assertion that the colour-swap is an involution.
\end{proposition}

\begin{proposition}[Names implemented by swaps; the dual calculus]\label{prop:dual}
Define the black calculus's reflection by $\quo{\,}_{\!\BS}:=\ktb\circ\swp\colon
\BS\to\BA$ and $\drp{\,}_{\!\BS}:=\swp\circ\ktb^{-1}$. Then by the coherence
$\swp\kt=\ktb\swp$,
\[
  \quo{\,}_{\!\BS} \;=\; \swp\circ\quo{\,}\circ\swp,
\]
so the two colours' quotes are swap-conjugate: one reflection determines the
other. The black calculus has black sets as processes and red sets as names; by
\cite{knot}, Corollary~2 the two interlock, each calculus's names being the
other's processes. A single reflective universe carries two dual rho calculi at
once.
\end{proposition}

\begin{remark}[The stratification is the grading]\label{rem:grading}
\cite{knot}, \S6.5 calls for a graded modality separating colour-respecting
(equivariant) computation from reflective computation, the price of breaking a
name's opacity. The four sorts are that grading. A computation that stays in
$\{\RS,\RA\}$ and uses only $\quo{\,},\drp{\,}$ relabels red processes and names
without ever observing black structure: this is the equivariant fragment, ordinary
nominal/$\pi$-style names with every guarantee intact. Applying $\swp$ to compute
with black set-operations on a name's interior is the reflective surplus, and the
type --- which sort one is in --- records it. Decidable structural identity of
names (\cite{knot}, \S6.5) is here the decidability of equality in $\RA$, settled
out of red's sight by black equality of the swapped processes.
\end{remark}

\section{The behaviour functor and the process universe}\label{sec:functor}

The knotted universe is a model of nominal sets (\cite{knot}, \S6.1): the model
$M$ of its Theorem~1 carries the $\Sym(\mathcal A)$-action with every element
finitely supported and supports the name-abstraction
$\nset X=(\Name\times X)/{\sim}$ with $\supp(\abs a x)=\supp(x)\setminus\{a\}$
(\cite{knot}, \S6.2). Write $\mathcal U$ for this nominal universe of red sets and
identify $\Name=\RA$ with $\Proc=\RS$ along the reflection of
Section~\ref{sec:swap}.

\begin{definition}[Behaviour functor]\label{def:functor}
Let $\B\colon\mathcal U\to\mathcal U$,
\[
  \B(R) \;=\; \mathcal P\!\big(\,\DK\times R\,\big),
\]
the (powerset of) context-label/residual pairs of Definition~\ref{def:ctx-lts},
with $\DK$ the Fire context-labels (Remark~\ref{rem:dT}). Because the input and
output labels carry an unrestricted message or continuation, $\DK$ is
class-indexed and $\B$ is taken with the full powerclass $\mathcal P$ rather than
$\Pfin$; its final coalgebra accordingly lives in the class-sized framework of
\cite{knot}, \S8.1, where it coincides with the colour-blind non-well-founded
universe of \cite{knot}, \S7. The \emph{commitment refinement} $\B^{\flat}(R)=
\Pfin\big(\Name\times\nset R\ \uplus\ \Name\times\Name\times R\ \uplus\ R\big)$
(input abstractions, output concretions, silent moves) presents the same
bisimilarity finitely (Remark~\ref{rem:commit}), assembled from constants,
products, coproducts, finite powerset, and the abstraction functor $\nset{(-)}$,
each preserving finite support and weak pullbacks; we use $\B^{\flat}$ when finite
branching matters.
\end{definition}

A $\B$-coalgebra $(R,\beta)$ equips every $r$ with its set of context-labelled
transitions --- the data of Section~\ref{sec:lts}. A coalgebra morphism is a
finitely-supported map commuting with the structure maps: a functional context
bisimulation.

\begin{definition}[The process universe]\label{def:proc}
$\Proc:=\nu\B$, the final $\B$-coalgebra in $\mathcal U$, with structure map
$\langle\mathsf{out},\mathsf{in},\tau\rangle\colon\Proc\to\B(\Proc)$, and
$\Name:=\swp(\Proc)$ presented as $\RA$ via $\kt$.
\end{definition}

Unfolded, an element of $\nu\B$ over a state is the set of that state's
transitions, the set of its residuals' transitions, and so on: the Aczel reading
$\{(a,P'):P\lts a P'\}$ with $\Name$-valued labels and an abstraction summand for
input. A loop is non-well-founded descent under the colour-blind $\in_\star$ of
\cite{knot}, \S7.2 --- a red process contains a black name whose red unfolding
contains a red process, without end. The construction's global
non-well-foundedness \emph{is} the recursion of process behaviour, available for
free where Aczel postulates it.

\section{A Milner-style denotation}\label{sec:denotation}

We now give $\sem{-}$ compositionally, following \cite{polyadic}. After Milner, a
guarded process commits to an \emph{action} carrying an \emph{agent}: an input
carries an \emph{abstraction} $\abs y P$ (it expects a name), an output a
\emph{concretion} $\conc Q$ (it offers one). Communication is application.

\begin{definition}[Commitment denotation]\label{def:milner}
Define $\sem{-}$ on processes by
\[
  \sem{\Zero}=\varnothing,\qquad
  \sem{\inp y x P}=\big\{\,\langle \sem{x}^{?},\ \abs y{\sem P}\rangle\,\big\},\qquad
  \sem{\out x Q}=\big\{\,\langle \sem{x}^{!},\ \conc{\sem Q}\rangle\,\big\},
\]
\[
  \sem{\drp x}=\drp{\sem x}=\swp(\kt^{-1}\sem x),\qquad
  \sem{P\Par Q}=\sem P \rpar \sem Q,
\]
and on names $\sem{\quo P}=\quo{\sem P}=\kt(\swp\,\sem P)\in\RA$. The tags
$\sem x^{?}$ and $\sem x^{!}$ abbreviate the context labels $\hole\Par\out x{(-)}$
and $\hole\Par\inp{(-)}x{(-)}$ of Definition~\ref{def:ctx-lts}, in their
finite commitment presentation (Remark~\ref{rem:commit}): the abstraction
$\abs y{\sem P}\in\nset{\nu\B}$ binds the received name, packaging the
class-indexed family of input context-transitions, and the concretion
$\conc{\sem Q}$ carries the message $\sem Q\in\nu\B$. Each guarded prefix denotes a
\emph{singleton} commitment, one cut at which $P$ meets its environment.
\end{definition}

\begin{definition}[Reaction; the swap fires]\label{def:react}
On commitments, application is
\[
  \abs y R \cdot \conc M \;=\; R\{\,\quo M\,/\,y\,\}
                        \;=\; R\{\,\kt(\swp\,M)\,/\,y\,\},
\]
and reaction between dual commitments on equivalent names is
\[
  \inferrule*[right=\textsc{React}]
    {\langle a^{?},\abs y R\rangle\in\sem P \\
     \langle a'^{!},\conc M\rangle\in\sem Q \\ a\nameq a'}
    {\sem{P}\Par\sem{Q} \ \react\ R\{\quo M/y\}\,\rpar\,(\text{residuals})} .
\]
The message $M$, a red set (a process), is sealed by $\quo{\,}=\kt\circ\swp$ into
the name $\quo M$ bound by the receiver: \emph{communication is precisely where the
colour-swap of Section~\ref{sec:swap} fires}. This is the set-theoretic content of
rho's \textsc{Comm}: the receiver learns not the process sent but its name.
\end{definition}

\begin{definition}[Semantic parallel]\label{def:rpar}
$\rpar\colon\nu\B\times\nu\B\to\nu\B$ is defined corecursively (by finality) as the
expansion
\[
  \begin{aligned}
  R\rpar R' \;=\;
   &\{\langle\alpha,\,T\rpar R'\rangle : \langle\alpha,T\rangle\in R\}
   \ \cup\
   \{\langle\alpha,\,R\rpar T'\rangle : \langle\alpha,T'\rangle\in R'\} \\
   &\cup\
   \{\langle\tau,\,U\rpar U'\rangle :
      \langle a^{?},\abs y U\rangle\in R,\ \langle a'^{!},\conc M\rangle\in R',\
      a\nameq a',\ U=\text{(redex residual)}\}
  \end{aligned}
\]
together with the symmetric $\tau$-set (a $\,?$-commitment of $R'$ with a
$\,!$-concretion of $R$). Here the silent reactions apply \textsc{React}
(Definition~\ref{def:react}), so $U$ is $\,U_0\{\quo M/y\}$ for the abstraction
body $U_0$. On immediate commitments $\rpar$ restricts to union, recovering the Milner intuition $\sem{P\Par Q}\supseteq\sem P\cup\sem Q$;
but the residuals carry the surviving context $T\rpar R'$, so $\rpar$ is not bare
union. This is what keeps $\out a 0\Par\out a 0\not\bis\out a 0$: the first, after
one $\out a{}$, retains the second output, where bare union would discard it.
\end{definition}

\begin{proposition}[The two denotations agree]\label{prop:agree}
The compositional $\sem{-}$ of Definitions~\ref{def:milner}--\ref{def:rpar}
coincides with the unique $\B^{\flat}$-coalgebra morphism from the syntactic
coalgebra $(\mathcal T/{\scong},\partial)$ of Section~\ref{sec:lts} into
$\nu{\B^{\flat}}$, where $\partial$ records each state's context-transitions in
their commitment presentation (Remark~\ref{rem:commit}).
\end{proposition}
\begin{proof}[Proof sketch]
Both are $\B^{\flat}$-coalgebra morphisms out of $(\mathcal T/{\scong},\partial)$:
the clauses of Definition~\ref{def:milner} read off exactly the commitments
$\partial$ assigns (singleton for each prefix, $\varnothing$ for $\Zero$), and
$\rpar$ is the corecursive map realising \textsc{Par} and \textsc{Comm}/\textsc{React}.
By finality the morphism is unique, so the two agree.
\end{proof}

\begin{theorem}[Adequacy and the reflection law]\label{thm:adequacy}
For all processes $P,Q$ and names $x$:
\begin{enumerate}[label=\textup{(\roman*)},nosep]
  \item $\sem{\quo P}=\quo{\sem P}$, $\sem{\drp x}=\drp x$, and
        $\sem{\drp{\quo P}}=\sem P$: the reflection law holds in the model,
        as $\swp\circ\swp=\mathrm{id}$ (Proposition~\ref{prop:reflection});
  \item $\sem{-}$ is a homomorphism: $\sem\Zero=\varnothing$,
        $\sem{P\Par Q}=\sem P\rpar\sem Q$, and the prefix clauses of
        Definition~\ref{def:milner}; in particular $P\scong Q\Rightarrow\sem P=\sem Q$,
        the commutative-monoid laws and reflection being special cases.
\end{enumerate}
\end{theorem}

\section{Full abstraction}\label{sec:fa}

We prove that $\sem{-}$ identifies exactly the context-bisimilar processes. The
argument is the standard coalgebraic one over the commitment refinement
$\B^{\flat}$, which finitely presents context bisimilarity; we then read the result
back through OSLF adequacy and note that, unlike the action-labelled case,
congruence is already in hand.

\begin{lemma}[Context $=$ coalgebraic bisimilarity]\label{lem:oc}
Context bisimilarity (Definition~\ref{def:bis}) coincides with $\B^{\flat}$-coalgebraic
bisimilarity. For the relational lifting $\overline{\B^{\flat}}$, a relation
$\mathcal R$ on $\mathcal T/{\scong}$ is a coalgebra bisimulation iff it is a
context bisimulation.
\end{lemma}
\begin{proof}[Proof sketch]
By Remark~\ref{rem:commit} the commitment presentation packages each class-indexed
context-transition family as a single abstraction or concretion, with no loss of
discriminating information: two states pass the same context labels iff their
commitment sets match. $\overline{\B^{\flat}}$ relates behaviours when their
output and silent summands match with $\mathcal R$-related residuals and their
input summands match in $\nset{\mathcal R}$; unfolding $\nset{\mathcal R}$ by its
defining equivalence ($\abs y R\mathrel{\nset{\mathcal R}}\abs y{R'}$ iff
$R\{a/y\}\mathrel{\mathcal R}R'\{a/y\}$ for all names $a$) reproduces the matching
of the offer-input context family, and the concretion case the offer-output family.
Hence the two notions of bisimulation coincide.
\end{proof}

\begin{lemma}[$\B^{\flat}$ is well behaved]\label{lem:wp}
$\B^{\flat}$ preserves weak pullbacks. Hence on any $\B^{\flat}$-coalgebra
$(C,\gamma)$, coalgebraic bisimilarity is the kernel of the unique morphism
$!_C\colon C\to\nu{\B^{\flat}}$: $c\mathrel{\sim_{\B^{\flat}}}d \iff
{!_C}(c)={!_C}(d)$.
\end{lemma}
\begin{proof}[Proof sketch]
$\Pfin$, products, coproducts and constants preserve weak pullbacks, and so does
the nominal abstraction functor $\nset{(-)}$ (its quotient by the freshness
equivalence is by a functional, hence pullback-stable, relation); weak-pullback
preservation is closed under these constructions, so $\B^{\flat}$ preserves them.
For a weak-pullback-preserving functor with a final coalgebra, behavioural
equivalence (the kernel of $!_C$) coincides with coalgebraic bisimilarity --- the
standard fact of universal coalgebra \cite{rutten}. Existence of $\nu{\B^{\flat}}$,
and its agreement with the class-sized $\nu\B$ on equality, is
Obligation~\ref{ob:exists}.
\end{proof}

\begin{theorem}[Full abstraction for context bisimilarity]\label{thm:fa}
For all rho processes $P,Q$,
\[
  \sem P=\sem Q \quad\Longleftrightarrow\quad P\bis Q .
\]
Moreover, since $\bis$ is a congruence (Proposition~\ref{prop:cong}), this is full
abstraction with respect to a congruence: $\sem{-}$ identifies $P$ and $Q$ exactly
when no context can tell them apart.
\end{theorem}
\begin{proof}
By Proposition~\ref{prop:agree}, $\sem{-}$ is the unique morphism
$!\colon(\mathcal T/{\scong},\partial)\to\nu{\B^{\flat}}$, so $\sem P=\sem Q$ iff
${!}(P)={!}(Q)$. By Lemma~\ref{lem:wp} on the syntactic coalgebra, ${!}(P)={!}(Q)$
iff $P,Q$ are coalgebraically bisimilar; by Lemma~\ref{lem:oc} that is context
bisimilarity $\bis$. Chaining, $\sem P=\sem Q\iff P\bis Q$. Congruence is
Proposition~\ref{prop:cong}.
\end{proof}

\begin{corollary}[Adequacy through OSLF]\label{cor:oslf}
Let $\mathcal L$ be the Hennessy--Milner logic the OSLF construction generates from
the context-labelled transition system, with modalities $\langle F\rangle\varphi$
indexed by the labels $\DK$ (Remark~\ref{rem:dT}). By OSLF adequacy
\cite{meredithwells}, $P\bis Q$ iff $P$ and $Q$ satisfy the same formulae of
$\mathcal L$. Combining with Theorem~\ref{thm:fa},
\[
  \sem P=\sem Q
  \ \Longleftrightarrow\
  P\bis Q
  \ \Longleftrightarrow\
  \forall\varphi\in\mathcal L.\ (P\models\varphi\Leftrightarrow Q\models\varphi).
\]
Equality in the model, context bisimilarity, and modal indistinguishability are one
relation seen three ways --- the Firing, Placing, and Witnessing readings of the
single label set $\DK$.
\end{corollary}

\begin{corollary}[Internal full abstraction; Aczel as a fragment]\label{cor:internal}
On $\nu{\B^{\flat}}=\Proc$ itself, two elements are context-bisimilar iff equal:
behavioural equivalence is the equality of the model. Aczel's strong extensionality
for CCS is recovered as the special case where the labels $\DK$ degenerate to an
opaque set of action names over which no further context structure is observed ---
the equivariant fragment of Section~\ref{sec:ledger}.
\end{corollary}

\paragraph{What the context labels bought.} An action-labelled presentation would
here have to negotiate congruence as an afterthought: strong bisimilarity in
name-passing calculi may fail to be preserved by input prefix unless one closes
under name instantiation, replacing $\bis$ by the largest bisimulation-congruence
$\bis^{c}$. Labelling by contexts removes the negotiation. The bisimulation game of
Definition~\ref{def:bis} already ranges over every environment, so
Proposition~\ref{prop:cong} delivers congruence directly (granted the relative
pushouts of Obligation~\ref{ob:rpo}), and Theorem~\ref{thm:fa} is full abstraction
for a congruence with no closure step. Rho's structure makes the relative-pushout
hypothesis especially plausible: there is no $\nu$, so no scope-extrusion
interaction with the cut, and a name is structured data, so name-for-name
substitution is set-theoretic substitution under $\kt\circ\swp$ with no fresh-name
side conditions.

\section{The risk ledger}\label{sec:ledger}

The accounting of \cite{knot} transfers, and locates where rho's recursion is paid
for.

\paragraph{Finite processes: the hereditarily-finite core, no infinite risk.}
A process with a finite synchronisation tree is a hereditarily-finite red set with
finitely many black-set atoms (its free names), living in the finitary core of
\cite{knot}, \S5.5, where the final coalgebra of the finite-powerset functor is
literal and exists at set size. Foundation and finite support are theorems, not
axioms; the commitment is the black empty set and one successor, no more.

\paragraph{Recursion: one $\omega$, black infinity, no primitive names.}
Recursion enters rho through reflection (Section~\ref{sec:rho}): a finite term can
loop by quoting and dereferencing its own continuation. Such a process is a
non-well-founded inhabitant of $\nu\B$, available only in the infinitary reading,
which by \cite{knot}, \S5.6 and \S8.1 costs exactly one $\omega$ --- black
infinity --- and lives, for the full powerset, in the class-sized final-coalgebra
framework. As with the completed datatype $\Lambda$ of \cite{knot}, \S6.2, each
finite-state recursive process is set-sized while the universe of all processes is
the one-$\omega$ object. No primitive name is ever reintroduced; names remain black
sets built from $\emptyset$.

\paragraph{Aczel's CCS model as the equivariant shadow.}
Stay within one colour, never apply $\swp$ to a name's interior, treat names as
opaque swap-only atoms (the equivariant fragment, \cite{knot}, \S6.5). The input
summand's abstraction degenerates, $\B$ reduces to $\B_0(R)=\Pfin(\mathrm{Act}\times R)$,
and $\nu{\B_0}$ is Aczel's CCS model (Corollary~\ref{cor:internal}). Aczel's
construction is the colour-respecting fragment; the rho surplus --- structured,
computed, content-addressed names, and the licence to cross the colour barrier
with $\drp{\,}$ --- is the reflective remainder, graded by the four sorts of
Section~\ref{sec:swap}.

\section{Toward a rigorous construction}\label{sec:obligations}

In the manner of \cite{knot}, \S8, the obligations a fully rigorous development
would discharge. For the finitary core each reduces to a statement about
hereditarily-finite sets, and several are corollaries of the model built.

\begin{obligation}[Existence of the final coalgebra, and $\nu\B^{\flat}\!\bis\nu\B$]\label{ob:exists}
Exhibit the final coalgebras via the two-sorted endofunctor
$H(R,B)=(B+\mathcal P_s R,\,R+\mathcal P_s B)$ of \cite{knot}, \S8.1, taken at the
\emph{final} coalgebra (greatest fixed point). For the commitment refinement
$\B^{\flat}$, finite branching makes it accessible and $\nu{\B^{\flat}}$ exists in
the nominal universe, at set size on the finitary core (the limit of the terminal
$\omega^{\mathrm{op}}$-chain). For the full context functor
$\B=\mathcal P(\DK\times-)$, the final coalgebra is the class-sized, colour-blind
non-well-founded universe of \cite{knot}, \S7. Show the two carry the same equality
--- that the commitment packaging of Remark~\ref{rem:commit} is a behavioural
equivalence --- so bisimilarity computed in $\nu{\B^{\flat}}$ is context
bisimilarity in $\nu\B$.
\end{obligation}

\begin{obligation}[The stratified signature and its coherence]\label{ob:strat}
Fix the four-sorted signature of Section~\ref{sec:swap} --- the sorts
$\RS,\RA,\BS,\BA$, the knots $\kt,\ktb$, the swap $\swp$ --- as an extension of the
$\Set[X]$ signature of \cite{knot}, and verify the equational theory: $\swp$ an
involution and a homomorphism for all set operations, $\swp\kt=\ktb\swp$, and
$\kt,\ktb$ isomorphisms. Then $\quo{\,},\drp{\,}$ are derived and
Proposition~\ref{prop:reflection} is a calculation. The collapse of \cite{knot}
--- $\kt,\ktb$ literal identities --- is recovered as a quotient of this theory,
and one checks the reflection law survives the quotient.
\end{obligation}

\begin{obligation}[Relative pushouts for rho's contexts]\label{ob:rpo}
Discharge the hypothesis of Proposition~\ref{prop:cong}: that the category of rho
contexts, with plugging as composition, possesses relative pushouts, so the
idem-pushout labels of Definition~\ref{def:ctx-lts} are well defined and the
Leifer--Milner congruence theorem applies. This is the GSLT minimal-context
construction \cite{meredithwells}; the cut lying always at a parallel composition
(no reaction under a prefix), and names being quoted structure rather than
$\nu$-bound, are the features that should make the relative pushouts exist and be
computable. With this discharged, $\bis$ is a congruence outright, and the
action-labelled closure-under-substitution question does not arise.
\end{obligation}

\begin{obligation}[Context correspondence and OSLF adequacy]\label{ob:oc}
Promote Lemma~\ref{lem:oc} to a theorem: define the relational lifting
$\overline{\B^{\flat}}$ precisely, including the abstraction summand, and prove
context bisimilarity equals $\overline{\B^{\flat}}$-bisimilarity. Relate $\bis$ to
reduction-barbed congruence so that Theorem~\ref{thm:fa} is full abstraction for
the intended equivalence; rho's lack of $\nu$ removes the closure under restriction
this normally requires. Finally, instantiate the OSLF construction on the
context-labelled system and prove the adequacy used in Corollary~\ref{cor:oslf}, so
that the modalities $\langle F\rangle\varphi$ indexed by $\DK$ characterise $\bis$.
\end{obligation}

\begin{obligation}[Substitution and the swap-firing reaction]\label{ob:subst}
Prove $\sem{P\{\quo M/y\}}=\sem P\{\quo{\sem M}/y\}$, i.e.\ that $\sem{-}$ is a map
of nominal algebras for substitution and that the application of
Definition~\ref{def:react} computes it, with the seal $\kt\circ\swp$ applied to the
message. This is the lemma underwriting \textsc{React} and clause (ii) of
Theorem~\ref{thm:adequacy} for input, carried by the abstraction functor's
evaluation $\nset R\times\Name\to R$.
\end{obligation}

\begin{obligation}[Barbed-congruence calibration]\label{ob:cong}
With congruence secured by Obligation~\ref{ob:rpo}, calibrate which behavioural
equivalence the model realises: confirm that context bisimilarity $\bis$ coincides
with strong barbed congruence for rho, so that Theorem~\ref{thm:fa} is full
abstraction for the equivalence practitioners use. Rho's structured, $\nu$-free
names are again the reason to expect the calibration to be tight where in $\pi$ it
is delicate.
\end{obligation}

\begin{obligation}[Placing the construction]\label{ob:place}
Locate the model among Aczel's CCS model \cite{aczel} (the equivariant fragment),
the presheaf and nominal coalgebraic models of name-passing
(Fiore--Turi \cite{fioreturi}, Stark \cite{stark}), of which this is the
reflective, $\nu$-free instance with $\Name\cong\Proc$, and the bialgebraic
semantics of Turi--Plotkin \cite{turiplotkin}, into whose abstract-GSOS format the
rho rules fit. The defining feature absent from all three --- $\Name\cong\Proc$ via
$\swp$ --- is the second knot.
\end{obligation}

\section{Conclusion}

Rho's defining move --- names are quoted processes, $\drp{\quo P}\scong P$ --- is
the colour-swap of the knotted universe, read through a stratification into four
sorts: $\quo{\,}=\kt\circ\swp$, $\drp{\,}=\swp\circ\kt^{-1}$, and the reflection
law is the involutivity $\swp\circ\swp=\mathrm{id}$, true by construction. The
dynamics follow Milner: a process is its set of commitments, input an abstraction
and output a concretion, communication their application --- and the application is
where the swap fires, sealing the sent process into the received name. The
denotation lands in the final coalgebra of a behaviour functor over the knotted
nominal universe, where the colour-blind non-well-foundedness supplies recursion
for free and finality makes bisimilar processes equal. Following the GSLT
discipline we labelled transitions by contexts --- the minimal environments in
which a process completes a reaction, which are at once the program/environment
cuts, the OSLF modality indices, and the locations in term structure --- and
expressed bisimulation over them; on that base we proved full abstraction,
$\sem P=\sem Q\iff P\bis Q$. The context labelling pays a structural dividend the
action-labelled account could not: the equivalence is a congruence by construction,
so full abstraction holds for a congruence with no closure step, granted only that
rho's contexts admit relative pushouts --- a hypothesis its $\nu$-free, quoted-name
structure makes especially plausible. The risk ledger is the source paper's: finite
processes cost nothing beyond the empty set, recursion costs one $\omega$, and not
one primitive name is ever spent. Aczel's model sits inside this one as its
equivariant shadow; the surplus is reflection, and reflection is the swap.

\subsection*{Acknowledgements}
This note develops Section~9.3 of \emph{A Knotted Universe}, with the stratified
reading of the colour barrier, the Milner-style commitment semantics over
context-labelled transitions, and the full-abstraction theorem as its additions.

\begin{thebibliography}{99}

\bibitem{aczel} Peter Aczel. \emph{Non-Well-Founded Sets}, volume 14 of CSLI
Lecture Notes. CSLI Publications, Stanford, CA, 1988.

\bibitem{knot} L. G. Meredith. A Knotted Universe: a new notion of reflective set
theories. Manuscript, 2026.

\bibitem{rho} L. Gregory Meredith and Matthias Radestock. A reflective
higher-order calculus. \emph{Electr.\ Notes Theor.\ Comput.\ Sci.},
141(5):49--67, 2005.

\bibitem{polyadic} Robin Milner. The Polyadic $\pi$-Calculus: A Tutorial. In
F.\,L. Bauer, W. Brauer, and H. Schwichtenberg, editors, \emph{Logic and Algebra
of Specification}, volume 94 of NATO ASI Series F, pages 203--246. Springer, 1993.
(Also Technical Report ECS-LFCS-91-180, University of Edinburgh, 1991.)

\bibitem{leifermilner} James J. Leifer and Robin Milner. Deriving bisimulation
congruences for reactive systems. In \emph{CONCUR 2000}, volume 1877 of Lecture
Notes in Computer Science, pages 243--258. Springer, 2000.

\bibitem{sewell} Peter Sewell. From rewrite rules to bisimulation congruences.
\emph{Theoret.\ Comput.\ Sci.}, 274(1--2):183--230, 2002.

\bibitem{meredithwells} L. G. Meredith, Christian Wells, et al. Minimal-context
transition systems and the observer-specified logic functor for graph-structured
lambda theories. Manuscript, forthcoming.

\bibitem{gp} Murdoch Gabbay and Andrew M. Pitts. A new approach to abstract syntax
with variable binding. \emph{Formal Asp.\ Comput.}, 13(3--5):341--363, 2002.

\bibitem{pitts} Andrew M. Pitts. \emph{Nominal Sets: Names and Symmetry in
Computer Science}, volume 57 of Cambridge Tracts in Theoretical Computer Science.
Cambridge University Press, 2013.

\bibitem{rutten} J.\,J.\,M.\,M. Rutten. Universal coalgebra: a theory of systems.
\emph{Theoret.\ Comput.\ Sci.}, 249(1):3--80, 2000.

\bibitem{milner} Robin Milner. \emph{Communication and Concurrency}. Prentice
Hall, 1989.

\bibitem{sw} Davide Sangiorgi and David Walker. \emph{The $\pi$-Calculus: A Theory
of Mobile Processes}. Cambridge University Press, 2001.

\bibitem{fioreturi} Marcelo Fiore and Daniele Turi. Semantics of name and value
passing. In \emph{16th Annual IEEE Symposium on Logic in Computer Science (LICS
2001)}, pages 93--104. IEEE, 2001.

\bibitem{stark} Ian Stark. A fully abstract domain model for the $\pi$-calculus.
In \emph{11th Annual IEEE Symposium on Logic in Computer Science (LICS 1996)},
pages 36--42. IEEE, 1996.

\bibitem{turiplotkin} Daniele Turi and Gordon Plotkin. Towards a mathematical
operational semantics. In \emph{12th Annual IEEE Symposium on Logic in Computer
Science (LICS 1997)}, pages 280--291. IEEE, 1997.

\bibitem{jmm} André Joyal and Ieke Moerdijk. \emph{Algebraic Set Theory}, volume
220 of London Mathematical Society Lecture Note Series. Cambridge University
Press, 1995.

\end{thebibliography}

\end{document}
